\documentclass[sigconf,anonymous]{aamas}
\settopmatter{printacmref=false}
\setcopyright{none}
\acmSubmissionID{2613}
\acmConference[AAMAS 2027]{Supplementary material for the 26th International Conference on Autonomous Agents and Multiagent Systems (AAMAS 2027)}{3--7 May 2027}{Hanoi, Vietnam}{}
\copyrightyear{2027}
\acmYear{2027}
\acmDOI{}
\acmISBN{}
\submissionType{Supplementary Material}
\usepackage{booktabs,amsmath,graphicx,tikz,hyperref}
\usetikzlibrary{arrows.meta,positioning,calc,fit,backgrounds}
\hypersetup{pdftitle={Supplement: When to Intervene},pdfauthor={}}
\raggedbottom
\emergencystretch=2em
\newsavebox{\tablebox}
\newcommand{\fitTable}[1]{\sbox{\tablebox}{#1}\ifdim\wd\tablebox>\linewidth\resizebox{\linewidth}{!}{\usebox{\tablebox}}\else\usebox{\tablebox}\fi}
\title{Supplementary Analysis: When to Intervene}
\author{Anonymous Author(s)}
\begin{document}
\onecolumn
\pagestyle{plain}
\begin{center}
{\fontsize{14}{16}\selectfont\bfseries Supplementary Analysis: When to Intervene}\par\medskip
Anonymous submission 2613
\end{center}
\section{Outcomes and paired comparisons}
The analyses use all 24 participants and 216 rounds (72 per condition). Every participant completed three rounds per condition, and all 24 contributed final-survey ratings. C denotes control, K constant assistance, and A adaptive assistance.

The four main outcomes are completion, correct pieces, adapted workload, and frustration. For each outcome, two-sided paired $t$ tests compare participants' three-round condition averages across K--C, A--C, and A--K. Holm correction jointly covers these twelve tests. A separate nineteen-test group covers three contrasts each for mental demand, physical demand, temporal demand, performance, and effort, and A--K contrasts for four assistance-experience items. These analyses and comparison groups are exploratory, not preregistered. Recorded round duration and final-survey ratings are descriptive.

Correct pieces range from 0 to 7, using the admissible orientation with the highest correct count. Completion corresponds to a solved seven-piece arrangement. Workload averages six seven-point burden-aligned items, reversing performance as $8-P$, then rescales the average $B$ to $(B-1)100/6$. Frustration is both a main outcome and a component of this composite; the two effects are not independent replications. Completion and correct pieces are also linked.

Table~\ref{tab:paired} reports pointwise confidence intervals, not simultaneous multiplicity-adjusted intervals. $p_{\mathrm{H}}$ denotes the Holm-adjusted value. Cohen's $d_z$ divides the paired mean difference by the standard deviation of participant differences. Significance is evaluated using $p_{\mathrm{H}}<.05$, not whether an unadjusted interval excludes zero.
\begin{table}[htbp]
\centering
\small
\caption{Main paired contrasts ($n=24$); completion differences are percentage points. Twelve-test Holm group.}
\label{tab:paired}
\fitTable{%
\begin{tabular}{llrlrr}
\toprule
Outcome & Contrast & Difference & 95\% CI & $d_z$ & $p_{\mathrm{H}}$ \\
\midrule
Completion & K--C & 44.44 & [32.23, 56.66] & 1.54 & $<.001$ \\
Completion & A--C & 34.72 & [21.94, 47.50] & 1.15 & $<.001$ \\
Completion & A--K & -9.72 & [-23.16, 3.71] & -0.31 & 0.444 \\
Correct pieces & K--C & 2.25 & [1.61, 2.89] & 1.47 & $<.001$ \\
Correct pieces & A--C & 2.00 & [1.46, 2.54] & 1.56 & $<.001$ \\
Correct pieces & A--K & -0.25 & [-0.71, 0.21] & -0.23 & 0.532 \\
Workload & K--C & -12.96 & [-17.40, -8.52] & -1.23 & $<.001$ \\
Workload & A--C & -16.40 & [-21.95, -10.84] & -1.25 & $<.001$ \\
Workload & A--K & -3.43 & [-9.66, 2.80] & -0.23 & 0.532 \\
Frustration & K--C & -0.76 & [-1.28, -0.25] & -0.63 & 0.027 \\
Frustration & A--C & -1.18 & [-1.70, -0.66] & -0.96 & $<.001$ \\
Frustration & A--K & -0.42 & [-0.90, 0.07] & -0.36 & 0.363 \\
\bottomrule
\end{tabular}}
\end{table}


\section{Puzzle- and period-adjusted robustness}
All adjusted models include condition, categorical puzzle identity, and categorical block period. Completion uses a binomial-logit generalized estimating equation with independent working correlation and participant-clustered Mancl--DeRouen bias-reduced covariance. The working correlation does not treat a person's rounds as independent. Correct pieces, workload, and numeric ratings use participant-fixed-effect ordinary least squares with CR1 participant-clustered covariance. CR1 applies the finite-sample factor $G/(G-1)\,(N-1)/(N-k)$, where $G$ is the participant count, $N$ the round count, and $k$ the parameter count. Contrasts use a $t$ reference with 23 degrees of freedom.

Completion contrasts are odds ratios; their tests use log-odds coefficients. Other adjusted effects are mean-score differences. They are not interchangeable with the observed paired differences. The twelve/nineteen correction groups remain the same. The adjusted models support the same main contrast conclusions: both assistance conditions outperform control on completion and pieces and reduce workload and frustration; none establishes an A--K difference. Constant--control frustration has adjusted $p_{\mathrm{H}}=.045$, compared with $.027$ in the paired analysis.
\begin{table}[htbp]
\centering
\small
\caption{Puzzle/block-adjusted main contrasts. Completion uses odds ratios (OR); other effects are score differences. Twelve-test Holm group.}
\label{tab:adjusted}
\fitTable{%
\begin{tabular}{llrlr}
\toprule
Outcome & Contrast & Effect & 95\% CI & $p_{\mathrm{H}}$ \\
\midrule
Completion (OR) & K--C & 8.16 & [3.84, 17.33] & $<.001$ \\
Completion (OR) & A--C & 5.43 & [2.41, 12.24] & 0.002 \\
Completion (OR) & A--K & 0.67 & [0.34, 1.30] & 0.802 \\
Correct pieces & K--C & 2.14 & [1.54, 2.74] & $<.001$ \\
Correct pieces & A--C & 1.95 & [1.41, 2.48] & $<.001$ \\
Correct pieces & A--K & -0.19 & [-0.74, 0.35] & 0.802 \\
Workload & K--C & -12.15 & [-16.94, -7.37] & $<.001$ \\
Workload & A--C & -15.04 & [-20.76, -9.31] & $<.001$ \\
Workload & A--K & -2.88 & [-9.38, 3.61] & 0.802 \\
Frustration & K--C & -0.72 & [-1.24, -0.20] & 0.045 \\
Frustration & A--C & -1.04 & [-1.58, -0.50] & 0.003 \\
Frustration & A--K & -0.32 & [-0.83, 0.18] & 0.802 \\
\bottomrule
\end{tabular}}
\end{table}


These are supplementary robustness checks, not a validation of the HR-rise cue or evidence of equivalence. With 24 participant clusters, covariance corrections and the $t$ reference remain approximations. Completion models converged with full-rank designs; base fitted probabilities ranged from .053 to .805. Numeric fixed-effect models can predict beyond bounded scales (14 piece predictions outside 0--7; one workload prediction outside 0--100); predictions were not clipped. A further categorical within-block-position check retained the direction of the main contrasts; its unadjusted diagnostic $p$ values are reported separately and do not replace the main tests.
\begin{table}[htbp]
\centering
\small
\caption{Additional within-block-position sensitivity: unadjusted diagnostic $p$ values, not replacements for family-adjusted tests.}
\label{tab:position}
\fitTable{%
\begin{tabular}{lrrr}
\toprule
Outcome & K--C & A--C & A--K \\
\midrule
Completion & $<.001$ & $<.001$ & 0.202 \\
Correct pieces & $<.001$ & $<.001$ & 0.449 \\
Workload & $<.001$ & $<.001$ & 0.353 \\
Frustration & 0.008 & $<.001$ & 0.181 \\
\bottomrule
\end{tabular}}
\end{table}


\section{Secondary ratings and study-wide experience}
Table~\ref{tab:secondary} provides every secondary contrast under the nineteen-test correction. Assisted-experience comparisons use 144 assisted rounds for adjusted models and all 24 participant averages for paired tests. Final helpfulness, efficacy, trust, and reliance summarize the study as a whole, not separate policies. Their means (SDs) are 5.00 (1.53), 4.83 (1.55), 4.79 (1.59), and 4.88 (1.48), respectively.
\begin{table}[htbp]
\centering
\small
\caption{Secondary ratings: nineteen-test Holm group, separate from the twelve main contrasts.}
\label{tab:secondary}
\fitTable{%
\begin{tabular}{llrlrr}
\toprule
Rating & Contrast & Difference & 95\% CI & Paired $p_{\mathrm{H}}$ & Model $p_{\mathrm{H}}$ \\
\midrule
Helpfulness & A--K & 0.25 & [-0.16, 0.66] & 1.000 & 1.000 \\
Timing & A--K & 0.56 & [0.16, 0.95] & 0.103 & 0.210 \\
Seamlessness & A--K & 0.08 & [-0.45, 0.61] & 1.000 & 1.000 \\
Frustration relief & A--K & 0.40 & [-0.02, 0.83] & 0.623 & 0.906 \\
Mental demand & K--C & -0.89 & [-1.24, -0.54] & $<.001$ & 0.001 \\
Mental demand & A--C & -1.08 & [-1.50, -0.67] & $<.001$ & $<.001$ \\
Mental demand & A--K & -0.19 & [-0.70, 0.32] & 1.000 & 1.000 \\
Physical demand & K--C & -0.38 & [-0.70, -0.05] & 0.320 & 0.717 \\
Physical demand & A--C & -0.60 & [-1.13, -0.07] & 0.320 & 0.736 \\
Physical demand & A--K & -0.22 & [-0.69, 0.24] & 1.000 & 1.000 \\
Temporal demand & K--C & -0.29 & [-0.64, 0.06] & 0.867 & 1.000 \\
Temporal demand & A--C & -0.36 & [-0.84, 0.12] & 1.000 & 1.000 \\
Temporal demand & A--K & -0.07 & [-0.66, 0.52] & 1.000 & 1.000 \\
Performance & K--C & 1.61 & [1.05, 2.17] & $<.001$ & $<.001$ \\
Performance & A--C & 2.01 & [1.42, 2.61] & $<.001$ & $<.001$ \\
Performance & A--K & 0.40 & [-0.18, 0.98] & 1.000 & 1.000 \\
Effort & K--C & -0.74 & [-1.08, -0.39] & 0.003 & 0.015 \\
Effort & A--C & -0.67 & [-1.11, -0.23] & 0.064 & 0.177 \\
Effort & A--K & 0.07 & [-0.27, 0.41] & 1.000 & 1.000 \\
\bottomrule
\end{tabular}}
\end{table}


\section{Condition orders and puzzle allocation}
Four participants followed each of CKA, CAK, KCA, KAC, ACK, and AKC. Every participant attempted all nine distinct puzzles. Order balance did not imply equal puzzle allocation by condition; Table~\ref{tab:allocation} shows the realized assignment.
\begin{figure}[htbp]
\centering
\begingroup
\fontsize{12}{13}\selectfont
\renewcommand{\small}{\fontsize{11}{12}\selectfont}
\resizebox{.75\linewidth}{!}{\definecolor{ink}{HTML}{1C1917}
\definecolor{muted}{HTML}{57534E}
\definecolor{rule}{HTML}{D6D3D1}
\definecolor{controlbg}{HTML}{F5F5F4}
\definecolor{adaptivebg}{HTML}{D1FAE5}
\definecolor{constantbg}{HTML}{FEF3C7}
\definecolor{adaptivefg}{HTML}{047857}
\definecolor{constantfg}{HTML}{B45309}
\definecolor{controlfg}{HTML}{44403C}

\newcommand{\condA}{\textcolor{adaptivefg}{\textbf{A}}}
\newcommand{\condK}{\textcolor{constantfg}{\textbf{K}}}
\newcommand{\condC}{\textcolor{controlfg}{\textbf{C}}}
\newcommand{\toord}{\,\textcolor{muted}{\textendash}\,}
\newcommand{\order}[3]{#1\toord#2\toord#3}

\begin{tikzpicture}[
  font=\sffamily,
  >={Stealth[length=2.4mm,width=1.8mm]},
  every path/.style={line cap=round,line join=round},
  edge/.style={line width=0.8pt,->},
  root/.style={
    draw=ink,line width=0.9pt,rounded corners=3.5pt,
    fill=ink,text=white,inner sep=9pt,
    font=\sffamily\bfseries,minimum height=11mm
  },
  cohort/.style={
    line width=0.8pt,rounded corners=3.5pt,
    inner sep=7pt,align=center,minimum width=22mm,
    font=\sffamily\bfseries
  },
  cond/.style={
    line width=0.9pt,rounded corners=3.5pt,
    inner sep=5pt,font=\sffamily\bfseries,minimum size=9.5mm
  },
  leaf/.style={
    line width=0.75pt,rounded corners=3.5pt,
    inner sep=7pt,align=center,minimum width=38mm
  },
  chip/.style={
    line width=0.7pt,rounded corners=2.5pt,
    inner sep=5pt,font=\sffamily,minimum height=7.2mm
  }
]

\node[root] (all) at (0,0) {24 participants};

\node[cohort,fill=controlbg,draw=controlfg,text=ink] (g1) at (4.0, 3.55) {Control first\\n=8};
\node[cohort,fill=adaptivebg,draw=adaptivefg,text=ink] (g2) at (4.0, 0) {Adaptive first\\n=8};
\node[cohort,fill=constantbg,draw=constantfg,text=ink] (g3) at (4.0,-3.55) {Constant first\\n=8};

\node[cond,fill=controlbg,draw=controlfg,text=controlfg] (c1) at (7.0, 3.55) {C};
\node[cond,fill=adaptivebg,draw=adaptivefg,text=adaptivefg] (c2) at (7.0, 0) {A};
\node[cond,fill=constantbg,draw=constantfg,text=constantfg] (c3) at (7.0,-3.55) {K};

\node[leaf,fill=controlbg,draw=controlfg] (l1) at (11.1, 4.35)
  {{\sffamily\small\color{muted}4 participants}\\[3pt]\order{\condC}{\condA}{\condK}};
\node[leaf,fill=controlbg,draw=controlfg] (l2) at (11.1, 2.75)
  {{\sffamily\small\color{muted}4 participants}\\[3pt]\order{\condC}{\condK}{\condA}};
\node[leaf,fill=adaptivebg,draw=adaptivefg] (l3) at (11.1, 0.80)
  {{\sffamily\small\color{muted}4 participants}\\[3pt]\order{\condA}{\condC}{\condK}};
\node[leaf,fill=adaptivebg,draw=adaptivefg] (l4) at (11.1,-0.80)
  {{\sffamily\small\color{muted}4 participants}\\[3pt]\order{\condA}{\condK}{\condC}};
\node[leaf,fill=constantbg,draw=constantfg] (l5) at (11.1,-2.75)
  {{\sffamily\small\color{muted}4 participants}\\[3pt]\order{\condK}{\condA}{\condC}};
\node[leaf,fill=constantbg,draw=constantfg] (l6) at (11.1,-4.35)
  {{\sffamily\small\color{muted}4 participants}\\[3pt]\order{\condK}{\condC}{\condA}};

\draw[edge,draw=controlfg] (all) -- (g1);
\draw[edge,draw=adaptivefg] (all) -- (g2);
\draw[edge,draw=constantfg] (all) -- (g3);
\draw[edge,draw=controlfg] (g1) -- (c1);
\draw[edge,draw=adaptivefg] (g2) -- (c2);
\draw[edge,draw=constantfg] (g3) -- (c3);
\draw[edge,draw=controlfg] (c1) -- (l1);
\draw[edge,draw=controlfg] (c1) -- (l2);
\draw[edge,draw=adaptivefg] (c2) -- (l3);
\draw[edge,draw=adaptivefg] (c2) -- (l4);
\draw[edge,draw=constantfg] (c3) -- (l5);
\draw[edge,draw=constantfg] (c3) -- (l6);

\draw[rule,line width=0.6pt] (0.15,-5.75) -- (13.15,-5.75);

\node[chip,fill=adaptivebg,draw=adaptivefg,text=ink] (kA) at (1.55,-6.45)
  {\condA\ \ Adaptive};
\node[chip,fill=constantbg,draw=constantfg,text=ink] (kI) at (4.15,-6.45)
  {\condK\ \ Constant};
\node[chip,fill=controlbg,draw=controlfg,text=ink] (kC) at (6.65,-6.45)
  {\condC\ \ Control};
\node[anchor=west,font=\sffamily\small,text=muted] at (8.05,-6.45)
  {Read left to right};

\end{tikzpicture}
}
\endgroup
\caption{Condition orders: four participants per order, 24 total.}
\end{figure}
\begin{table}[htbp]
\centering
\small
\caption{Puzzle allocation; each condition has 72 attempts.}
\label{tab:allocation}
\fitTable{%
\begin{tabular}{lrrrr}
\toprule
Puzzle & Control & Constant & Adaptive & Total \\
\midrule
1 & 8 & 11 & 5 & 24 \\
2 & 9 & 10 & 5 & 24 \\
3 & 9 & 12 & 3 & 24 \\
4 & 8 & 7 & 9 & 24 \\
5 & 5 & 8 & 11 & 24 \\
6 & 6 & 3 & 15 & 24 \\
7 & 6 & 8 & 10 & 24 \\
8 & 11 & 7 & 6 & 24 \\
9 & 10 & 6 & 8 & 24 \\
\bottomrule
\end{tabular}}
\end{table}


\section{Reproducibility and interpretation}
The analysis used Python 3.12.14, NumPy 2.3.3, SciPy 1.16.2, pandas 2.3.3, statsmodels 0.14.5, and patsy 1.0.1. Independent checks reproduced paired tests, cluster-covariance calculations, and Holm adjustment. No rounds were excluded or imputed in these analyses. Aggregate comparison CSVs accompany this supplement; raw participant files and recordings are not included.

The policies share an available preset-hint list and retrieval rule, not a demonstrated equal intervention dose. Effects therefore concern the assistance packages, not timing alone or robot retrieval in isolation. Recorded duration reflects round end-marking for solved and unsolved attempts and is not treated as verified active solving time. Nonsignificance is not evidence that two policies are equivalent.
\end{document}
